\scanpage{164}\mspage{22}
% archive sheet number: 164
$\underline U,\underline V$ \deleted{\unsure{à} \unsure{Ob}}\unsure{de} $S_{/x}$, si $\underline U,\underline V$ non vides (ce qui signifie que les sous-objets de $x$ dans $\unsure{\mathscr X|\mathcal X}^{\wedge}$ qu'ils définissent sont ``non vide'', \unsure{vu|ou} que les objets de $S$ sont ``non vides'' \unsure{ds} dans $\mathscr X^{\wedge}$) alors $\underline U\cap\underline V$ est non vide. \unsure{OBS|Obs} que $\underline U,\underline V$ est engendré \inserted{chacun} par une seule flèche $a\overset u\longrightarrow x$, $b\overset v\longrightarrow a$. Le \deleted{\unsure{un}} crible intersection est formé des objets $c\longrightarrow x$ de $S_{/x}$ tels que $c\longrightarrow x$ se factorise \deleted{par} en
\diagram{\begin{tikzcd}[ampersand replacement=\&,row sep=small,column sep=small]
\&a\\
c\arrow[ur,dashed]\arrow[dr,dashed]\&\\
\&b
\end{tikzcd}}
, donc \deleted{l'exi} dire que ce crible n'est pas vide, signifie \deleted{que les flèches} que $\mathrm{PS}1$ est satisfait pour les deux flèches données $u,v$. Donc c) équivaut : la condition $\mathrm{PS}1$ $\unsure{e|\in}$.

Ainsi on a établi
\[
b'\Longrightarrow b\Longrightarrow a\Longrightarrow(c\Longleftrightarrow d\Longleftrightarrow e),
\]
reste à prouver que $e\Longrightarrow b'$. Soit donc
\diagram{\begin{tikzcd}[ampersand replacement=\&,row sep=small,column sep=small]
U\arrow[dr]\&\&V\arrow[dl]\\
\&W\&
\end{tikzcd}}
un diagramme d'objets ``non-vides'' \deleted{avec} avec $W$ $0$-connexe, \deleted{\unsure{il a que} \illegible{2} que le produit} \inserted{il faut prouver} \inserted{que le produit} fibré est non vide. \deleted{Comme il y a un objet au bas, on a là \unsure{cet}} Considérant $a\longrightarrow U$ et $b\longrightarrow V$ avec $a,b\in S$, on est ramené au cas où $U=a$, $V=b$ \deleted{dans} dans $S$.

Prouvons que si $a\times_W b$ est ``vide'', alors $W$ est déconnexe. Pour ceci \deleted{fixons} \deleted{et considérons}\inserted{considérons} la sous-catégorie $C$ de $S_{/W}$ formée des $x\longrightarrow W$ tels que $a\times_W x$ soit ``non vide'', \deleted{ou que le facteur \unsure{soit}} Désignons par $w(a),w(x)$ les ``images'' de $a,x$ dans $W$, cela signifie que $w(a)\cap w(x)\ne\varnothing$.

\deleted{Comme $a$ est irréductible, on voit \unsure{de suite}}
